\documentclass[10pt, a4paper]{article}
\usepackage[T2A]{fontenc}
\usepackage[utf8]{inputenc}
\usepackage{enumitem}
\usepackage[russian]{babel}
\usepackage{amsmath, amsfonts, amssymb, amsthm, tikz}
\usepackage{geometry}
\geometry{top=1.5cm, bottom=1.5cm, left=1.5cm, right=1.5cm}


\title{\vspace{-1cm}\textbf{Конспект по дискретной математике}\vspace{-0.5cm}}
\author{}
\date{}

\def\divby{%
  \mathrel{\text{\vbox{\baselineskip.65ex\lineskiplimit0pt\hbox{.}\hbox{.}\hbox{.}}}}%
}

\begin{document}
\maketitle
\vspace{-0.8cm}

\section*{Список используемых обозначений}
\begin{tabular}{ll}
    $\lfloor x \rfloor$ & --- Округление числа $x$ вниз \\
    $\lceil x \rceil$ & --- Округление числа $x$ вверх \\
    $\iff$ & --- Тогда и только тогда, когда \\

\end{tabular}

\vspace{0.5cm}
\hrule

\section*{\centerline{\Large\textbf{\S \,\, Методы доказательств}}}
\vspace{0.4cm}

\subsection*{Конструктивное док-во (пример)}

\textbf{Пример:} Докажите, что пифагоровых троек бесконечно много.
\[ 3^2 + 4^2 = 5^2 \]
\[ (3k)^2 + (4k)^2 = (5k)^2 \text{ --- бесконечная серия, } \forall k \in \mathbb{N} \]

\subsection*{Контрпример}

\textbf{Вопрос:} Верно ли, что все слова русского языка не содержат более 4 согласных подряд? Нет, «контрпример» имеет 5.

\subsection*{Прямое док-во}

Если $a \divby c$ и $b \divby d$, то $ab \divby cd$.
\[ \left.
\begin{aligned}
a \divby c &\implies \exists k: a = ck \\
b \divby d &\implies \exists l: b = dl
\end{aligned}
\right\} \implies ab = ckdl \implies ab \divby cd \text{ по опр.} \]

\subsection*{Док-во от противного}
\noindent \hfill \[ A \longrightarrow B \iff \overline{B} \longrightarrow \overline{A} \]
\noindent
\textbf{Пример:} Доказать, что простых чисел бесконечно много.
\begin{proof}[Док-во]
\noindent
Пусть не так, тогда $p_1, \dots, p_n$ --- все простые числа. \\
Рассмотрим $x = p_1 \cdot p_2 \cdot p_3 \cdot \dots \cdot p_n + 1$, заметим, что 
$\left.
\begin{aligned}
x > p_i, \; \forall i \\
x \ndivs p_i, \; \forall i
\end{aligned}
\right\} \implies x \text{ --- новое простое число.}$
\end{proof}

\subsection*{Оценка + пример}

\textbf{Вопрос:} Какое наибольшее кол-во не бьющих друг друга ладей?

\begin{proof}[Док-во]
Пример на 8:
\begin{center}
\begin{picture}(80,80)
\multiput(0,0)(0,10){9}{\line(1,0){80}}
\multiput(0,0)(10,0){9}{\line(0,1){80}}
\multiput(5,5)(10,10){8}{\circle*{4}}
\end{picture}
\end{center}

\noindent От противного: Пусть 9, тогда 2 будут на одной горизонтали и будут бить друг друга. \\
\textit{На самом деле здесь мы пользуемся принципом Дирихле.}
\end{proof}

\subsection*{Принцип Дирихле}

\textbf{Утв.} Даны $n$ клеток и $n+1$ кроликов, то существует клетка с $\ge 2$ кроликами.

\paragraph{Утв. Обобщенный принцип Дирихле} 
Даны $n$ клеток, $k$ кроликов, то:
\begin{enumerate}
    \item $\exists$ клетка с $\ge \lceil \frac{k}{n} \rceil$ кроликами.
    \item $\exists$ клетка с $\le \lfloor \frac{k}{n} \rfloor$ кроликами.
\end{enumerate}

\begin{proof}[Док-во обобщенного принципа Дирихле:]
\noindent 
\begin{enumerate}
    \item От противного: пусть в каждой $< \lceil \frac{k}{n} \rceil$. В клетках $a_1, \dots, a_n$ кроликов. \\
    \[ a_i < \left\lceil \frac{k}{n} \right\rceil \implies a_i \le \left\lceil \frac{k}{n} \right\rceil - 1 < \frac{k}{n} \]
    \[ k = \sum a_i < n \cdot \frac{k}{n} = k \quad ?! \]

    \item Пусть в каждой $> \lfloor \frac{k}{n} \rfloor$.
    \[ a_i > \left \lfloor \frac{k}{n} \right \rfloor \implies a_i \ge \left \lfloor \frac{k}{n} \right \rfloor + 1 > \frac{k}{n} \]
    \[ k = \sum a_i > n \cdot \frac{k}{n} = k \quad ?! \]
\end{enumerate}
\end{proof}

\textbf{Пример:} 40 человек. Докажите, что есть месяц с $\ge 4$ днями рождения и есть месяц с $\le 3$. \\

\begin{proof}{\textit{по принципу Дирихле:}}
\begin{enumerate}
    \item Пусть в каждом $\le 3 \implies$ всего дней рожд $\le 36$, но людей $40$?!
    
    \item Пусть в каждом $\ge 4 \implies$ всего $\ge 48$, но людей $40$?!
\end{enumerate}
\end{proof}

\paragraph{Утв. (Бесконечный принцип Дирихле)} 
Есть бесконечное число кроликов и конечное число клеток, тогда $\exists$ клетка, в которой бесконечное число кроликов.
\begin{proof}[Док-во:]
Пусть нет, тогда в каждой клетке конечное число кроликов, но $\text{конечное} \times \text{конечное} = \text{конечное} < \text{бесконечное}$.
\end{proof}

\subsection*{Комбинаторные док-ва}

Вместо счета объектов и проверки их равенства, можно установить соответствие.

\subsection*{Подсчет двумя способами}
Часто используется в док-ве от противного: «что-то одновременно больше и меньше нуля».

\subsection*{Принцип крайнего}

Хорошая штука, когда не хватает условий. Часто работает с методом от противного, но это не обязательно.
\noindent        
\textbf{Задача:} Докажите, что среди $n\ge 3$ точек на плоскости есть такая, что все углы с другими точками $<180^\circ$. 
\begin{proof}[Док-во:]
    Рассмотрим самую левую точку и из таких самую нижнюю, заметим, что она подойдет, так как нет точек левее ее и $180^\circ$ тоже не бывает. 
\end{proof}

\subsection*{Инвариант --- «что-то, что не меняется»}

\textbf{Задача:} На доске написано число $a$. Разрешается его стереть и написать либо $3a$, либо $a+8$, либо $\frac{a}{5}$, если $a \divs 5$. Можно ли из $1$ получить $2$?

\noindent \textit{Док-во:} Нет, т.к. число всегда нечётное. Инвариант --- чётность.

\noindent \textbf{Вопрос:} А из числа $2$ в $4$? Нет, но инвариант сложнее.

\subsection*{Вспомогательная раскраска}

\textbf{Задача:} Из шахматной доски вырезали 2 противоположных угла. Можно ли её разрезать на доминошки $1 \times 2$?

\noindent \textit{Док-во:} Отметим шахматную раскраску. Пусть отрезанные углы чёрные, тогда у нас 32 белых, 30 чёрных, а каждая доминошка содержит по одной клетке каждого цвета $\implies$ не выйдет.

\noindent Инвариант: кол-во белых $-\,\, \text{кол-во}$ чёрных.

\subsection*{Полуинвариант --- «что-то, что меняется лишь в одну сторону»}

\textbf{Пример:} предыдущая задача, но можно ещё вырезать чёрные клетки как действие.

\noindent Полуинвариант: кол-во белых $-\,\, \text{кол-во}$ чёрных. Он только увеличивается, изначально 2, а должен стать 0 ?!

\subsection*{Индукция}

Нужна, чтобы доказывать утверждение для всех $n$ ($\forall n$).\\

\noindent \textbf{Пример 1:} $1 + 2 + 3 + \dots + n = \frac{n(n+1)}{2}$

\noindent База: $n = 1 \implies 1 = \frac{1(1+1)}{2}$

\noindent Переход: $n = k \implies n = k+1$. Хотим для $k+1$:
\[ 1 + 2 + 3 + \dots + k + (k+1) = \underbrace{\frac{k(k+1)}{2}}_{\text{= по ИП}} + (k+1) = \frac{(k+1)(k+2)}{2} = \frac{(k+1)((k+1)+1)}{2} \quad \checkmark \]

\subsection*{Возвратная индукция --- переход из $1, \dots, n \implies n+1$}

\textbf{Пример:} Алгоритм Евклида заканчивается. Делаем индукцию по $k = n + m$, переход от всех меньших k к k.

\noindent Переход: $(n, m) = (n - m, m)$ --- для второй пары заканчивается по ИП (индукционному предположению), т.к. $(n - m) + m < n + m$.

\subsection*{Переход: $n \implies n+x$}

\textbf{Задача:} Какую сумму можно выдать монетами номинала 3 и 5? \\
Докажем, что любую сумму $n \ge 8$ можно выдать.

\noindent Переход: $n \implies n+3$. Выдадим 3, а $n$ --- по ИП.

\noindent База: 
$n = 8 = 3 + 5; \quad n = 9 = 3 + 3 + 3; \quad n = 10 = 5 + 5$.

\subsection*{Переход: $2^n \implies 2^{n+1}$ и $x \implies x-1$}

\textbf{Задача:} Доказать, что любое число раскладывать в сумму степеней двойки, где каждая не более одного раза.

\noindent Переход: $n \implies 2n$ и $n \implies 2n+1$.

\section*{\centerline{\Large\textbf{\S \,\, Введение в теорию множеств}}}
\vspace{0.4cm}

Множество — первичное математическое понятие, которому не дано строгое математическое определение. Множество состоит из элементов. Множество задается набором своих элементов.

\vspace{0.2cm}
\noindent \textbf{Обозн.} $A$ --- множество, $x$ --- элемент, $x \in A$.

\vspace{0.4cm}

\noindent \textbf{Опр.}
Множество $A$ называется подмножеством множества $B$, если $\forall x \in A \implies x \in B$.

\vspace{0.2cm}
\noindent \textbf{Обозн.} $A \subset B$.

\vspace{0.2cm}
\noindent \textbf{Замеч.} Иногда используют символ $\subseteq$, что означает нестрогое включение, т.е. множество \(A\) является подмножеством \(B\) и может быть равным ему (\(A = B\)). Если $A \neq B$, то $A$ называется собственным подмножеством и обозначается $A \subsetneq B$.

\vspace{0.4cm}

\noindent \textbf{Утв.}
Если $A \subset B$ и $B \subset A$, то $A = B$.

\vspace{0.2cm}
\begin{proof}[Док-во:]
$\forall x \in A \implies x \in B \quad \text{и} \quad \forall x \in B \implies x \in A$
\end{proof}

\noindent \textbf{Опр.}
Множество, не содержащее ни одного элемента, называется пустым.

\vspace{0.2cm}
\noindent \textbf{Обозн.} $\varnothing$

\vspace{0.6cm}

\noindent \textbf{Замеч. (Парадокс Рассела / Брадобрея)}

\vspace{0.2cm}
\noindent Рассмотрим множество всех множеств, не содержащих самих себя. 
Вопрос: содержит ли это множество само себя?

\[ A = \{ B \mid B \notin B \} \]
\newpage
\noindent Разберем два возможных случая:
\begin{itemize}
    \item Если $A \in A$, то по определению множества $A \implies A \notin A$.
    \item Если $A \notin A$, то по определению множества $A \implies A \in A$ \quad ?!
\end{itemize}

На самом деле, нельзя так задавать множества. Нужно писать:
\[ A = \{ B \in \Omega \mid \text{условие на } B \} \]

\noindent Поэтому везде будем считать, что все множества лежат в некотором большом множестве $\Omega$.

\vspace{0.3cm}

\subsection*{Как задавать множества?}

\begin{enumerate}
    \item Поэлементно, например $A = \{3, \text{стул}, \text{мем с котиком}\}$
    \item Через условия, которые называются предикатами. \\
    $A = \{ x \in \mathbb{N} \mid x \;\divby\; 5 \}$ \\
    $P(x) = x \;\divby\; 5$ --- предикат.
    
    \textit{Замечание.} Не любое условие является предикатом, далее мы будем это игнорировать
\end{enumerate}

\vspace{0.3cm}
\noindent \textbf{Пример (Парадокс Берри)}

\vspace{0.2cm}
\noindent Рассмотрим множество $A = \{ x \in \mathbb{N} \mid x \text{ можно описать менее, чем 12 словами русского языка} \}$.

\vspace{0.2cm}
\noindent С одной стороны, множество $A$ конечно. Тогда существует наименьшее число $y$, которое в нем не лежит. Но тогда: $y$ --- это «наименьшее число, которое нельзя описать менее, чем двенадцатью словами русского языка» $\implies y \in A$ \quad ?!

\vspace{0.3cm}
\subsection*{Операции над множествами}

\begin{itemize}
    \item Множество $A \cap B = \{ x \mid x \in A \land x \in B \}$ называется \textbf{пересечением} ($\land$ --- конъюнкция --- логическая связка «и»).
    \item Множество $A \cup B = \{ x \mid x \in A \lor x \in B \}$ называется \textbf{объединением} ($\lor$ --- дизъюнкция --- логическая связка «или»).
\end{itemize}

\vspace{0.2cm}
\noindent \textit{Замечание.} Формально нужно писать $\{ x \in \Omega \mid \dots \}$, но в наивной теории множеств это подразумевается.

\begin{itemize}
    \item Множество $\overline{A} = \{ x\mid  x \in \Omega \land x \notin A \}$ называется \textbf{дополнением}
    \item Множество $A \setminus B = \{ x \mid x \in A \land x \notin B \}$ называется \textbf{разностью}
    \item Множество $A \Delta B = (A \setminus B) \cup (B \setminus A)$ называется \textbf{симметрической разностью} \\
    \small (множество элементов, лежащих ровно в одном из множеств)
\end{itemize}

\vspace{0.2cm}
\subsection*{Свойства операций над множествами:}

\begin{itemize}
    \item \textbf{Коммутативность:} 
    $A \cap B = B \cap A, \quad A \cup B = B \cup A, \quad A \Delta B = B \Delta A$
    
    \vspace{0.2cm}
    \begin{proof}[Док-во на примере объединения:]
    \[ A \cup B = \{ x \in \Omega \mid x \in A \lor x \in B \} \]
    \[ B \cup A = \{ x \in \Omega \mid x \in B \lor x \in A \} \]
    $x \in A \lor x \in B \iff x \in B \lor x \in A$ (т.к. логическая дизъюнкция коммутативна)
    \end{proof}

    \vspace{0.1cm}
    \item \textbf{Ассоциативность:}
    $A \cap (B \cap C) = (A \cap B) \cap C$ \\
    $A \cup (B \cup C) = (A \cup B) \cup C$ \\
    $A \Delta (B \Delta C) = (A \Delta B) \Delta C$
    \begin{proof}[Док-во:]
    \[ A \Delta B = \{ x \mid x \in A \text{ xor } x \in B \} \]
    Логическое исключающее «или» (xor) ассоциативно, откуда следует ассоциативность для любого числа множеств.
    \end{proof}

    \vspace{0.3cm}
    \item \textbf{Дистрибутивность:} \\
    $A \cup (B \cap C) = (A \cup B) \cap (A \cup C)$ \\
    $A \cap (B \cup C) = (A \cap B) \cup (A \cap C)$
    
    \vspace{0.2cm}
    \begin{proof}[Док-во:]
    \[ A \cup (B \cap C) = \{ x \mid x \in A \lor (x \in B \land x \in C) \} \]
    \[ x \in A \lor (x \in B \land x \in C) \iff (x \in A \lor x \in B) \land (x \in A \lor x \in C) \]
    Полученное выражение в точности совпадает с определением правой части:
    \[ (A \cup B) \cap (A \cup C) = \{ x \mid (x \in A \lor x \in B) \land (x \in A \lor x \in C) \} \]
    \end{proof}
    \item \textbf{Закон двойного отрицания:}
    $\overline{\overline{A}} = A$

    \item \textbf{Законы де Моргана:} $\overline{A \cap B} = \overline{A} \cup \overline{B}$, \quad $\overline{A \cup B} = \overline{A} \cap \overline{B}$
    \end{itemize}
    
\noindent Множество упорядоченных пар $A \times B = \{ (x, y) \mid x \in A \text{ и } y \in B \}$ называется \textbf{декартовым произведением}

\vspace{0.4cm}

\noindent \textbf{Опр.}
\textbf{(Бинарным) отношением} $R$ называется любое подмножество декартова произведения $A \times B$.

\vspace{0.2cm}
\noindent \textbf{Пример:}
Отношение нестрогого порядка $\le$ для натуральных чисел: \\
${\le} = \{ (1, 1), (1, 2), (1, 3), (2, 2), (2, 3), \dots \}$ 

\vspace{0.2cm}
\noindent Вместо записи $(x, y) \in R$ пишут: $xRy$.

\vspace{0.4cm}
\noindent \textbf{Опр.}
\textbf{Функция} --- это такое отношение $R$, которое обладает следующими свойствами:
\[ \begin{cases} 
    \forall x \in A \;\; \exists y \in B : (x, y) \in R \\ 
    \forall x \in A \;\; \forall y_1, y_2 \in B : \big( (x, y_1) \in R \land (x, y_2) \in R \big) \implies y_1 = y_2 
\end{cases} \]

\vspace{0.4cm}

\noindent \textbf{Опр.}
Два множества называются \textbf{равномощными}, если между ними существует взаимно однозначное соответствие (\textbf{биекция}).

\vspace{0.4cm}

\noindent \textbf{Опр.}
Множество $A$ \textbf{не больше по мощности}, чем множество $B$, если существует соответствие элементов $B$ элементам $A$ (но не обязательно всем элементам $B$), являющееся \textbf{инъекцией}.

\vspace{0.2cm}
\noindent \textbf{Обозн.} $|A| = |B| \iff A \sim B$ \quad $|A| \le |B| \iff A \preceq B$

\vspace{0.4cm}

\noindent \textbf{Пример:} Интуитивно кажется, что $|\mathbb{N}| \prec |\mathbb{Z}|$, но на самом деле их мощности равны: $|\mathbb{N}| = |\mathbb{Z}|$.

Сопоставим каждому целому числу натуральное по формуле 
\begin{equation*}
    f(x) = \begin{cases} 
2x+1, & \text{если } x \ge 0, \\ 
2|x|, & \text{если } x < 0. 
\end{cases}
\end{equation*}

Заметим, что оно взаимно однозначное.

\vspace{0.4cm}

\noindent \textbf{Опр.}
Множество называется \textbf{счетным}, если оно равномощно множеству натуральных чисел $\mathbb{N}$ (т.е. все его элементы можно пронумеровать).

\vspace{0.4cm}

\noindent $\mathbb{Q}$ --- счётное множество, $\frac{p}{q} \in \mathbb{Q}, \quad p \in \mathbb{Z}, \quad q \in \mathbb{N}$


\begin{center}
\begin{tikzpicture}[scale=1.2, >=stealth]
    \draw[->, thick] (-3.5,0) -- (3.5,0) node[right] {$p \in \mathbb{Z}$};
    \draw[->, thick] (0,-0.5) -- (0,4.5) node[above] {$q \in \mathbb{N}$};
    
    
    \foreach \p in {-3,-2,-1,0,1,2,3} {
        \foreach \q in {1,2,3} {
            \filldraw[black] (\p,\q) circle (1.5pt);
            \ifnum\p=0
                \node[below right, scale=0.65, text=gray] at (\p,\q) {$0$};
            \else
                \node[below right, scale=0.65, text=gray] at (\p,\q) {$\frac{\p}{\q}$};
            \fi
        }
    }
    
    \draw[->, red, thick] (0,1) -- (1,1);
    \draw[->, red, thick] (1,1) -- (0,2);
    \draw[->, red, thick] (0,2) -- (-1,1);
    \draw[->, red, thick] (-1,1) -- (-2,1);
    \draw[->, red, thick] (-2,1) -- (-1,2);
    \draw[->, red, thick] (-1,2) -- (0,3);
    \draw[->, red, thick] (0,3) -- (1,2);
    \draw[->, red, thick] (1,2) -- (2,1);
    \draw[->, red, thick] (2,1) -- (3,1);
    \draw[->, red, thick] (3,1) -- (2,2);
    \draw[->, red, thick] (2,2) -- (1,3);
    \draw[->, red, thick] (1,3) -- (0,4);
    \draw[->, red, thick] (0,4) -- (-1,3);
    \draw[->, red, thick] (-1,3) -- (-2,2);
    \draw[->, red, thick] (-2,2) -- (-3,1) node[left, scale=0.8, text=red] {$\dots$};

    % Метки на горизонтальной оси p
    \foreach \x in {-3,-2,-1,1,2,3} \node[below=1mm, scale=0.8] at (\x,0) {\x};
    \node[below left=1mm, scale=0.8] at (0,0) {0};
    
    % Метки на вертикальной оси q
    \foreach \y in {1,2,3,4} \node[left=1mm, scale=0.8] at (0,\y) {\y};
\end{tikzpicture}
\end{center}


\vspace{0.4cm}

\noindent \textbf{Теорема}
Отрезок $[0, 1]$ --- несчётное множество.

\vspace{0.3cm}

\begin{proof}[Доказательство (Диагональный метод Кантора):]
Пусть мы занумеровали все числа отрезка. Запишем их в десятичной записи:
\begin{align*}
    1. & \quad 0, \textbf{1}11111\dots \\
    2. & \quad 0,1\textbf{2}00\dots \\
    3. & \quad 0,10\textbf{1}001\dots \\
    4. & \quad 0,146\textbf{9}2\dots
\end{align*}

\noindent Построим число, которое отличается от первого в первом знаке после запятой, от второго --- во втором знаке после запятой и т.д. Будем брать любую цифру, кроме той, которая там стоит, и кроме цифры $9$, чтобы не возникло периодического переноса разряда в виде $0.999\dots = 1.0$.

\vspace{0.2cm}
\noindent Тогда это новое число отличается от всех чисел в списке. Следовательно, его в этом списке нет, что приводит к противоречию.
\end{proof}

\vspace{0.2cm}
\noindent \textbf{Опр.}
Мощность отрезка $[0, 1]$ называется \textbf{континуумом}.

\end{document}

